\chapter{Scientific Text}
\label{chap:bt-scientific}

\noindent
\begin{center}
\begin{tikzpicture}
\node[draw=ThesisNeutral, fill=ThesisPaper, rounded corners=3pt, line width=0.5pt,
      inner sep=10pt] {%
\begin{tabular}{@{}l !{\hspace{6pt}\textcolor{ThesisNeutral}{\vrule width 0.3pt}\hspace{6pt}} l@{}}
\begin{minipage}[t]{0.40\textwidth}
\footnotesize\color{ThesisInk}
{\tiny\sffamily\color{ThesisNeutral} FRANÇAIS}\\[5pt]
\textit{L'insuffisance mitrale (IM) ischémique est une valvulopathie fonctionnelle,
restrictive, liée au remodelage ventriculaire gauche post-infarctus et qui
complique environ 20\% des cardiopathies ischémiques. Sa présence grève lourdement
le pronostic de la maladie coronarienne, tant en termes de morbidité que de
mortalité. Encore mal connue, elle suscite depuis quelques années de nombreux
travaux visant à préciser ses mécanismes physiopathologiques, ses facteurs
pronostiques et sa prise en charge thérapeutique.}
\end{minipage}
&
\begin{minipage}[t]{0.40\textwidth}
\footnotesize\color{ThesisNeutral}
{\tiny\sffamily\color{ThesisNeutral} ENGLISH · TRANSLATION}\\[5pt]
\textit{Ischaemic mitral regurgitation is a functional, restrictive valve disease,
linked to post-infarction left-ventricular remodelling, which complicates about
20\% of ischaemic heart diseases. Its presence weighs heavily on the prognosis of
coronary disease, in terms of both morbidity and mortality. Still poorly understood,
it has prompted in recent years a large body of work seeking to clarify its
pathophysiological mechanisms, its prognostic factors and its therapeutic
management.}
\end{minipage}
\end{tabular}%
};
\end{tikzpicture}
\captionof{figure}{A research article about myocardial infarction; French original, English translation.}
\label{fig:bt-scientific}
\end{center}
\medskip

% === P1 — HOOK ===
\Cref{fig:bt-scientific} is about the same heart attack as the last chapter. This
time it covers one of its later complications. Yet it looks nothing like the note.
The sentences are complete. Each one follows from the one before. The passage even
stops to name what is still unknown and to promise it will study it. There is also
a lot of writing like this, and unlike the clinical note it is free to read. It is
tempting to treat it as the clinical language written with more care and easy to
find. It is not. The research article is abundant and public, but it is a language
of its own.

% === P2 ===
The two texts describe the same disease, so the article can look like a tidier
version of the note. The real difference is purpose. The rhetorician Carolyn Miller
asked why texts settle into the forms they do. She answered that a genre is a
recurring action in a community: to know a genre is to know what it is for and what
it does \citep{miller_genre_1984}. By that test the two are different genres because
they answer different needs. The clinical note exists to record a patient and guide
their care. The research article exists to put a claim before the field and open it
to challenge. This difference in purpose shapes how each one is written.

% === P3 ===
The article's purpose shows in its shape. The linguist John Swales studied how
researchers open their papers. He found a pattern so regular that he named it:
first stake out a territory, then carve a niche by showing what is missing from it,
then fill that niche with the present work \citep{swales_genre_1990}. The example
in \Cref{fig:bt-scientific} runs through all three. It claims a territory by noting
that ischaemic mitral regurgitation complicates about a fifth of ischaemic heart
disease and weighs on the prognosis. It opens a niche with \emph{still poorly
understood}. And it fills that niche by pointing to the body of work, its own
included, that sets out to explain it. The opening move is to argue that the paper
needs to exist. A clinical note never makes this argument. It has no gap to open and
nothing to justify. It simply records.

% === P4 ===
The body of the article follows the same logic. Its four parts, introduction,
methods, results, and discussion, retrace the steps of a controlled study: a
question, the steps used to answer it, what was observed, and what it means.
The arrangement feels timeless, but it is recent. One survey looked at 1{,}297
original articles from four leading medical journals between 1935 and 1985. Not one
used this structure in 1935. By the 1970s more than four in five did, and by the
1980s it was the only accepted form for an original paper \citep{sollaci_imrad_2004}.
A format only a few decades old now carries the shape of scientific reasoning in its
headings alone. A reader can jump to the question, the method, or the verdict
without reading the rest.

% === P5 ===
The clinic has a structure of its own, and it follows a different kind of reasoning.
The problem-oriented record was built to organise care around a patient's problems.
Its familiar subjective-objective-assessment-plan note follows the loop of diagnosis
and management rather than experiment \citep{weed_records_1968}. The article asks whether a
hypothesis holds. The note asks what is wrong with the patient and what to do next.
Both are about the same medicine, but they reason in different ways. One is not a
worse version of the other.

% === P6 ===
The difference reaches down to the sentence. An article does more than arrange its
argument. Sentence by sentence it keeps signalling how sure it is. Ken Hyland's
studies of academic writing call this running layer metadiscourse: the writer's own
handling of the claim, including the words that tune how strongly it is made
\citep{hyland_metadiscourse_2005, hyland_tse_2004}. Hedges such as \emph{may} or
\emph{suggests} hold a statement back. Boosters such as \emph{clearly} or
\emph{shows} push it forward. \emph{These results suggest that X may contribute to
Y} is a claim offered for debate. The plain \emph{X causes Y} would be far bolder.
The note hedges too, but rarely, and for another reason: its \emph{rule out
myocardial infarction} is an instruction to act.

% === P6b — Biber quantitative register ===
These impressions can be counted. The linguist Douglas Biber measured which
features cluster together across many spoken and written texts. He placed scientific
prose about as far from everyday conversation as written English gets: dense with
nouns, long words, and prepositions, and thin on the pronouns and verbs of speech
\citep{biber_variation_1988}. The density is not even the same within one article.
The methods read flat and impersonal, the discussion hedged and personal. So the
same shift in certainty that Hyland reads in the words shows up as a measurable
difference between sections \citep{biber_finegan_1994}.

% === P6c — phraseology / Nwogu ===
The pattern reaches the smallest scale. Kevin Nwogu catalogued the medical
article's rhetorical moves in detail, eleven of them across the usual sections
\citep{nwogu_medical_1997}. Its set phrases are just as fixed: frames such as
\emph{the role of} or \emph{has been shown to} recur with a narrow,
discipline-specific set of fillers \citep{luzon_collocational_2000}. From its
overall plan down to its smallest phrases, the article is built for a purpose the
clinical note does not share.

% === P7 ===
This hedging does real work. Watching how scientists write up their results, Bruno
Latour noticed that a finding hardens into a fact as its qualifiers fall away from
one citation to the next. Once enough people repeat it unchallenged, \emph{X may
contribute to Y} becomes simply \emph{X causes Y} \citep{latour_laboratory_1979}. The
hedges that fill the literature are therefore the visible trace of claims still
being argued over. A model that flattens them mistakes a provisional finding for a
settled fact.

% === P8 ===
There is a more formal way to say all of this. A sublanguage, in the sense the last
chapter borrowed from Zellig Harris, is the limited grammar a specialised domain
settles into, regular enough to be written down and handed to a computer. Carol
Friedman and colleagues built two such grammars, one from patient reports and one
from biology articles. They found the two were truly distinct: each domain needs its
own grammar to be processed \citep{friedman_two_2002}. The lesson for what follows
is plain. A model trained on the abundant public literature learns to argue about
medicine. It does not learn to record a patient. More articles, on their own, do not
bring us closer to the clinic.

% === P9 ===
So the scientific article is very different from the clinical note. It is public
where the note is private. It spells things out where the note leaves them unsaid.
It weighs each claim where the note simply acts. It is abundant and free to read,
but it is its own language. There is a third kind of medical text, more abundant
still and written for the widest audience of all, the patients themselves. Whether
that brings us any nearer to the clinic is where we turn next.
